% HoTT foundations: fixed official2013 public-release source; no HIT/cubical substitution.
\chapter{Paths, Equivalences and Homotopy Levels}
\label{sec:p07-hott-paths-and-equivalences}

An identity type carries identifications, not merely a yes-or-no answer to
an equality question. Identifications themselves have identity types.
This permits path operations and their higher laws without forcing every
type to be a set. The presentation follows Chapters~2 and~3 of the official
2013 \emph{HoTT Book} source at its June 19 public-release commit, with the
specific Chapter~4 and~7 dependencies identified below. A separately pinned
official errata companion corrects typed formulas; the rolling current
first-edition PDF is not silently substituted for that source
\parencite[ch.~2--3]{HoTT2013PublicReleaseFoundations}.

\section{The source theory and the host interface}
\label{sec:p07-hott-theory}

\begin{hbconvention}{def:p07-hott-core}{Identity, elimination and equality standing}
Write $x=_Ay$ for $\operatorname{Id}_A(x,y)$ and
$\mathsf{refl}_x:x=_Ax$. The source's $\equiv$ denotes judgmental equality;
$=$ in a type denotes an identity type. These are different typed roles.
The host's \hyperref[sheet:mltt-intensional]{MLTT sheet} supplies Id/J and
J-computation without equality reflection. A J-motive may depend on both
endpoints and the path, with its specified value on reflexivity.

For the following source presentation we also use the book's dependent
$\Sigma$ elimination with computation on pairs, its $\Pi$-\(\eta\)
uniqueness principle, and the unit, empty and two-element inductive types.
Those principles are not inferred from the host sheet's projections and
beta rules alone. Universes are predicative and cumulative, with levels
chosen to contain each displayed type. No UIP, proof irrelevance, global
set axiom, resizing, excluded middle or choice axiom is added silently.
Function extensionality, univalence and propositional truncation will be
declared separately where used.
\end{hbconvention}

Based path induction varies the other endpoint: for fixed $a$, a motive
over $(y,p:a=y)$ is determined from $(a,\mathsf{refl}_a)$. It does not
allow treating an arbitrary loop $p:a=a$ with both endpoints fixed as
reflexivity. The contractibility of the entire based-path total space,
proved below, is compatible with nontrivial loops in a particular fiber
\parencite[\S1.12; \S2.1]{HoTT2013PublicReleaseFoundations}.

\begin{hbdefinition}{def:p07-hott-path-operations}{Composition, inversion and loops}
For $p:x=y$ and $q:y=z$, write $p\cdot q:x=z$, traversing $p$ first.
Define composition by the source's double J construction, so only
$\mathsf{refl}_x\cdot\mathsf{refl}_x\equiv\mathsf{refl}_x$ is claimed
as its specified computation. Define $p^{-1}:y=x$ by J, with
$(\mathsf{refl}_x)^{-1}\equiv\mathsf{refl}_x$.
For a pointed type $(A,a)$, $\Omega(A,a)=(a=a,\mathsf{refl}_a)$;
iterating this construction defines its higher loop types.
\end{hbdefinition}

\begin{hbtheorem}{thm:p07-hott-path-laws}{proved=annex:p07-hott-path-laws}{Weak groupoid laws and cancellation}
There are paths witnessing left and right units, both inverse laws,
$(p^{-1})^{-1}=p$, and associativity
$p\cdot(q\cdot r)=(p\cdot q)\cdot r$.
Composition with a fixed path, on either side, has a quasi-inverse obtained
by composing with its inverse, so path cancellation is justified.
These laws are identities between paths, not judgmental equalities for
arbitrary inputs. The canonical J-built associator satisfies the pentagon
between the five reassociations of four composable paths.
\end{hbtheorem}

\section{Action, dependent action and transported comparisons}
\label{sec:p07-hott-transport}

\begin{hbdefinition}{def:p07-hott-transport}{Transport, ap and apd}
J defines
\[
 \mathsf{transport}^{P}(p,u):P(y),\quad
 \mathsf{ap}_f(p):f(x)=f(y),\quad
 \mathsf{apd}_s(p):\mathsf{transport}^{P}(p,s(x))=s(y),
\]
for $P:A\to\mathcal U$, $p:x=y$, $u:P(x)$, $f:A\to B$ and
$s:\Pi_{x:A}P(x)$. On reflexivity these compute respectively to
$u$, $\mathsf{refl}_{f(x)}$ and $\mathsf{refl}_{s(x)}$.
A dependent path from $u:P(x)$ to $v:P(y)$ over $p$ means
$\mathsf{transport}^{P}(p,u)=v$ in $P(y)$, not an ill-typed equality
between elements of different fibers.
\end{hbdefinition}

\begin{hbtheorem}{thm:p07-hott-transport-laws}{proved=annex:p07-hott-transport-laws}{All functorial and transport comparisons used here}
Ap preserves composition and inverses, commutes with composition of
functions and sends the identity function's action to the original path.
Transport obeys
\[
 \mathsf{transport}^{P}(p\cdot q,u)
 =\mathsf{transport}^{P}(q,\mathsf{transport}^{P}(p,u)),
\]
has inverse transport along $p^{-1}$, and satisfies
\[
 \mathsf{transport}^{P\circ f}(p,u)
 =\mathsf{transport}^{P}(\mathsf{ap}_f(p),u).
\]
For $k_x:P(x)\to Q(x)$ it also satisfies
$\mathsf{transport}^{Q}(p,k_x(u))=k_y(\mathsf{transport}^{P}(p,u))$.
Transport in a constant family is propositionally the identity; composing
that comparison with $\mathsf{ap}_f(p)$ gives $\mathsf{apd}_f(p)$.
In equality families the typed formulas are
\[
 \begin{aligned}
 \mathsf{transport}^{z\mapsto(a=z)}(p,r)&=r\cdot p,\\
 \mathsf{transport}^{z\mapsto(z=a)}(p,r)&=p^{-1}\cdot r,\\
 \mathsf{transport}^{z\mapsto(f(z)=g(z))}(p,r)
   &=\mathsf{ap}_f(p)^{-1}\cdot r\cdot\mathsf{ap}_g(p).
 \end{aligned}
\]
For dependent $f,g:\Pi_x B(x)$, the last formula instead uses
$\mathsf{apd}_f(p)^{-1}\cdot
\mathsf{ap}_{\mathsf{transport}^{B}(p,-)}(r)\cdot\mathsf{apd}_g(p)$.
\end{hbtheorem}

\begin{hbdefinition}{def:p07-hott-homotopy}{Homotopy is pointwise identification}
For sections $f,g:\Pi_xP(x)$, define
$f\sim g=\Pi_x(f(x)=g(x))$. Reflexivity, inversion and composition are
pointwise path operations. A homotopy is not automatically a path $f=g$.
\end{hbdefinition}

\begin{hbtheorem}{thm:p07-hott-homotopy-natural}{proved=annex:p07-hott-homotopy-natural}{Naturality and the self-composite law}
For $H:f\sim g$ between $A\to B$ and $p:x=y$,
\[
 H(x)\cdot\mathsf{ap}_g(p)=\mathsf{ap}_f(p)\cdot H(y).
\]
If $H:k\sim\mathsf{id}_A$, then
$H(k(x))=\mathsf{ap}_k(H(x))$. Both are higher paths, obtained without
assuming function extensionality.
\end{hbtheorem}

\begin{hbtheorem}{thm:p07-hott-eckmann-hilton}{proved=annex:p07-hott-eckmann-hilton}{Second-loop commutativity with typed whiskering}
For $\alpha,\beta:\mathsf{refl}_a=\mathsf{refl}_a$ in the path type
$a=a$, $\alpha\cdot\beta=\beta\cdot\alpha$.
This does not assert commutativity of arbitrary first loops.
\end{hbtheorem}

\section{Equivalence data and a coherence which must be proved}
\label{sec:p07-hott-equivalences}

\begin{hbdefinition}{def:p07-hott-equivalence}{The source's distinct equivalence notions}
For $f:A\to B$, a quasi-inverse is
$(g,\eta,\epsilon)$ with $g:B\to A$, $\eta:g f\sim\mathsf{id}_A$ and
$\epsilon:f g\sim\mathsf{id}_B$; its type is $\mathsf{qinv}(f)$.
The left and right inverse types are respectively
$\mathsf{linv}(f)=\Sigma_g(g f\sim\mathsf{id}_A)$ and
$\mathsf{rinv}(f)=\Sigma_g(f g\sim\mathsf{id}_B)$.
$\mathsf{biinv}(f)$ is their product, permitting different inverse functions.

A half-adjoint equivalence adds
$\tau:\Pi_{a:A}(\mathsf{ap}_f(\eta_a)=\epsilon_{f(a)})$ to the quasi-inverse
data. Its type is $\mathsf{ishae}(f)$. Following the source's final choice
in \S4.5, put $\mathsf{isequiv}(f)=\mathsf{ishae}(f)$ and
$A\simeq B=\Sigma_{f:A\to B}\mathsf{isequiv}(f)$.
This definition does not prove that raw quasi-inverse data have coherent
triangles, unique witnesses or contractible fibers; the following results do.
\end{hbdefinition}

\begin{hbtheorem}{thm:p07-hott-adjointification}{proved=annex:p07-hott-adjointification}{A constructed half-adjoint triangle}
Given $(g,\eta,\epsilon):\mathsf{qinv}(f)$, keep $g,\eta$ and replace
the other homotopy by the well-typed composite
\[
 \epsilon'_b=\epsilon_{f(g(b))}^{-1}\cdot
                     \mathsf{ap}_f(\eta_{g(b)})\cdot\epsilon_b.
\]
Then $\mathsf{ap}_f(\eta_a)=\epsilon'_{f(a)}$ is inhabited.
Half-adjoint data yield quasi-inverse data by forgetting this triangle.
Bi-invertible data also yield a quasi-inverse by explicitly joining the
possibly different left and right inverse functions.
\end{hbtheorem}

The original 2013 displayed concatenations in this construction were
reversed. The separately dated official errata confirms that error and
the triangle's orientation. The displayed construction above is typed in
the stated traversal convention, and the annex supplies its proof
\parencite[\S2.4; \S4.2]{HoTT2013PublicReleaseFoundations}
\parencite[entries for Theorem4.2.3]{HoTTOfficialErrata2023Foundations}.

\section{Pairs and functions have different path principles}
\label{sec:p07-hott-type-formers}

\begin{hbtheorem}{thm:p07-hott-sigma-paths}{proved=annex:p07-hott-sigma-paths}{Paths in total spaces, including both round trips}
For $u:P(x)$ and $v:P(y)$ there is a specified equivalence
\[
 ((x,u)=(y,v))\simeq
 \Sigma_{p:x=y}(\mathsf{transport}^{P}(p,u)=v).
\]
It computes at reflexivity to the pair of reflexivities; its inverse is
denoted $\mathsf{pairpath}$. It yields the propositional pair uniqueness
$w=(\pi_1w,\pi_2w)$ and a lift
$\mathsf{lift}(u,p):(x,u)=(y,\mathsf{transport}^{P}(p,u))$.
Transport in a product is componentwise. For
$F(x)=\Sigma_{u:P(x)}Q(x,u)$ it satisfies
\[
 \mathsf{transport}^{F}(p,(u,z))=
 (\mathsf{transport}^{P}(p,u),
       \mathsf{transport}^{Q}(\mathsf{lift}(u,p),z)).
\]
None of these says that equality of pairs with the same first coordinate
must arise over the reflexivity path of that coordinate.
\end{hbtheorem}

\begin{hbdefinition}{def:p07-hott-funext}{The explicitly assumed function-extensionality axiom}
For $f,g:\Pi_xP(x)$, J gives
$\mathsf{happly}:(f=g)\to(f\sim g)$. Where stated, assume this particular
map is an equivalence. A chosen inverse is $\mathsf{funext}$; its
pointwise beta comparison and path eta comparison are propositional.
They are not supplied by J or by judgmental function eta alone.
\end{hbdefinition}

\begin{hbtheorem}{thm:p07-hott-function-paths}{proved=annex:p07-hott-function-paths}{Function paths and transport under funext}
Assuming funext, $\mathsf{happly}(\mathsf{funext}(h),x)=h(x)$ and
$\mathsf{funext}(\mathsf{happly}(r))=r$. Function path inversion and
composition correspond to the pointwise operations.
For families $P,Q:X\to\mathcal U$, $p:x=y$ and $f:P(x)\to Q(x)$,
\[
 \mathsf{transport}^{z\mapsto(P(z)\to Q(z))}(p,f)
 =\lambda v.\mathsf{transport}^{Q}
            (p,f(\mathsf{transport}^{P}(p^{-1},v))).
\]
For $Q:\Sigma_{x:X}P(x)\to\mathcal U$, dependent transport obeys
the corresponding pointwise formula: transport the argument backwards,
apply $f$, and transport its result along the inverse of the lift of
$p^{-1}$. In particular the dependent-path comparison is
\[
 (\mathsf{transport}^{F}(p,f)=g)\simeq
 \Pi_{u:P(x)}\bigl(
 \mathsf{transport}^{Q}(\mathsf{lift}(u,p),f(u))
                  =g(\mathsf{transport}^{P}(p,u))\bigr),
\]
where $F(x)=\Pi_{u:P(x)}Q(x,u)$. The nondependent version replaces
$Q(x,u)$ by $Q(x)$.
\end{hbtheorem}

The distinction between mere existence of maps in both directions and a
specified equivalence with both inverse comparisons matters at function
types, just as it does at pair types
\parencite[\S2.6--2.9]{HoTT2013PublicReleaseFoundations}.

\section{Contractibility is coherent data, not connectedness}
\label{sec:p07-hott-contractibility}

\begin{hbdefinition}{def:p07-hott-contractibility}{The three bottom levels}
Define
\[
 \begin{aligned}
 \mathsf{isContr}(A)&=\Sigma_{a:A}\Pi_{x:A}(a=x),\\
 \mathsf{isProp}(A)&=\Pi_{x,y:A}(x=y),\\
 \mathsf{isSet}(A)&=\Pi_{x,y:A}\Pi_{p,q:x=y}(p=q).
 \end{aligned}
\]
A contraction has a specified center and a dependent family of paths.
It is stronger than a point together with merely inhabited path types.
A mere proposition may be empty; a contractible type cannot be empty.
A set is a type whose identity types are mere propositions, not a type
which must have only one point.
\end{hbdefinition}

\begin{hbtheorem}{thm:p07-hott-contractibility-laws}{proved=annex:p07-hott-contractibility-laws}{Canonical paths and removal of contractible data}
If $(a,c):\mathsf{isContr}(A)$, every $p:x=y$ satisfies
$p=c_x^{-1}\cdot c_y$; hence every path type in $A$ is contractible.
A type is contractible exactly when it is inhabited and a mere proposition,
or equivalently when it is equivalent to $\mathbf1$.
The based-path total space $\Sigma_{x:A}(a=x)$ is contractible even when
the loop fiber $a=a$ need not be.

Retracts of contractible types are contractible. If all $P(x)$ are
contractible, $\Sigma_xP(x)\simeq A$; if $A$ is contractible with center
$a$, $\Sigma_xP(x)\simeq P(a)$. A contractible base with contractible
fibers has contractible total space. Assuming funext, pointwise contractible
fibers also give a contractible dependent product, and
$\mathsf{isContr}(A)$ is a mere proposition.
\end{hbtheorem}

\begin{hbtheorem}{thm:p07-hott-prop-set}{proved=annex:p07-hott-prop-set}{Propositions, sets and proof standing}
Every mere proposition is a set, and a mere proposition's path types are
contractible. Logical equivalence of mere propositions yields type
equivalence by constructing inverse comparisons.
Assuming funext, $\mathsf{isProp}(A)$ and $\mathsf{isSet}(A)$ are themselves
mere propositions, and dependent products of mere propositions are mere
propositions. For a proposition-valued family $P$, paths in
$\Sigma_xP(x)$ are determined up to equivalence by their first projection.
Products of mere propositions remain mere propositions, whereas
$\mathbf1+\mathbf1$ does not; likewise $\Sigma_{x:\mathbf2}\mathbf1$
is not a mere proposition despite its proposition-valued fibers.
\end{hbtheorem}

\section{The actual fiber criterion and the bridges between notions}
\label{sec:p07-hott-fibers}

\begin{hbdefinition}{def:p07-hott-fiber}{Homotopy fibers and the map-level overload}
$\mathsf{fib}_f(b)=\Sigma_{a:A}(f(a)=b)$.
The source also writes
$\mathsf{isContr}(f)=\Pi_{b:B}\mathsf{isContr}(\mathsf{fib}_f(b))$.
Here the operand is a map, unlike $\mathsf{isContr}(A)$ for a type.
These source notations are retained with their types distinguished.
\end{hbdefinition}

\begin{hbtheorem}{thm:p07-hott-fiber-equivalence}{proved=annex:p07-hott-fiber-equivalence}{Centers, contractions and their coherent converse}
Paths from $(a,p)$ to $(a',p')$ in $\mathsf{fib}_f(b)$ are equivalent to
$\Sigma_{\gamma:a=a'}(\mathsf{ap}_f(\gamma)\cdot p'=p)$.
Half-adjoint data $(g,\eta,\epsilon,\tau)$ make each fiber contractible,
with center $(g(b),\epsilon_b)$ and contraction having first component
$\mathsf{ap}_g(p)^{-1}\cdot\eta_a$.
Conversely, specified fiber centers and contractions construct $g$,
$\epsilon$, $\eta$ and the actual triangle $\tau$.
Thus quasi-inverse data yield contractible fibers through the proved
adjointification construction, not through bare inverse uniqueness.
\end{hbtheorem}

\begin{hbtheorem}{thm:p07-hott-equivalence-bridges}{proved=annex:p07-hott-equivalence-bridges}{The three good notions are mere propositions and equivalent}
Assuming funext, $\mathsf{ishae}(f)$, $\mathsf{biinv}(f)$ and
$\mathsf{isContr}(f)$ are mere propositions and are pairwise equivalent.
For a map with a quasi-inverse, the individual types of left inverses
and right inverses are contractible. Equivalences admit identity,
inversion and composition, with actual composite inverse homotopies.
\end{hbtheorem}

\begin{hbtheorem}{thm:p07-hott-path-equivalence}{proved=annex:p07-hott-path-equivalence}{An equivalence induces equivalences on paths}
For $f:A\simeq B$, the map
$\mathsf{ap}_f:(a=a')\to(f(a)=f(a'))$ is an equivalence.
For half-adjoint inverse data its quasi-inverse sends $q$ to
$\eta_a^{-1}\cdot\mathsf{ap}_g(q)\cdot\eta_{a'}$.
Both round trips use naturality and the triangle; the second cannot be
justified by ignoring that coherence.
\end{hbtheorem}

These bridges follow the source's \S4.2--4.5. Its raw $\mathsf{qinv}$
type is deliberately not identified with a mere-propositional predicate.
The countercase below explains why its extra data cannot simply be erased
\parencite[\S2.4, \S2.11, \S4.1--4.5]{HoTT2013PublicReleaseFoundations}.

\section{Homotopy levels and their closure proofs}
\label{sec:p07-hott-levels}

\begin{hbdefinition}{def:p07-hott-levels}{The source's recursive n-type predicate}
For integers $n\geq-2$, define
\[
 \begin{aligned}
 \mathsf{is}\mbox{-}(-2)\mbox{-}\mathsf{type}(A)&=\mathsf{isContr}(A),\\
 \mathsf{is}\mbox{-}(n+1)\mbox{-}\mathsf{type}(A)
   &=\Pi_{x,y:A}\mathsf{is}\mbox{-}n\mbox{-}\mathsf{type}(x=y).
 \end{aligned}
\]
The index is the shifted natural-number induction of source \S7.1.
$(-1)$-types are mere propositions and $0$-types are sets.
This is a predicate on a given type, not the construction of an
$n$-truncation of every type.
\end{hbdefinition}

\begin{hbtheorem}{thm:p07-hott-level-closure}{proved=annex:p07-hott-level-closure}{Cumulativity, retracts, sums and products}
The hierarchy is upward closed. Retracts, hence equivalent types, preserve
each level. If the base and every fiber are $n$-types, their dependent sum
is an $n$-type. Assuming funext, a dependent product of $n$-types is an
$n$-type over any base, and each predicate
$\mathsf{is}\mbox{-}n\mbox{-}\mathsf{type}(A)$ is a mere proposition.
These claims include the contractible base case, not merely set closure.
\end{hbtheorem}

\section{The additional axioms and concrete countercases}
\label{sec:p07-hott-axioms-and-boundaries}

\begin{hbdefinition}{def:p07-hott-univalence}{A particular map, not an arbitrary identification}
J supplies $\mathsf{idtoeqv}:(A=_{\mathcal U}B)\to(A\simeq B)$ by transport
in the family $X\mapsto X$. For the results marked univalent, assume
that this particular map is an equivalence in the chosen universe.
Its chosen inverse is $\mathsf{ua}$. Univalence is not manufactured
from Id/J, and no judgmental computation for $\mathsf{ua}$ is asserted.
\end{hbdefinition}

\begin{hbtheorem}{thm:p07-hott-univalence-comparison}{proved=annex:p07-hott-univalence-comparison}{Propositional computation and a universe which is not a set}
Under univalence and the declared funext assumption,
$\mathsf{transport}^{X\mapsto X}(\mathsf{ua}(e),a)=e(a)$,
$\mathsf{ua}(\mathsf{id}_A)=\mathsf{refl}_A$,
$\mathsf{ua}(e)\cdot\mathsf{ua}(d)=\mathsf{ua}(d\circ e)$, and
$\mathsf{ua}(e)^{-1}=\mathsf{ua}(e^{-1})$.
The universe containing $\mathbf2$ is not a set: the Boolean interchange
equivalence gives a loop at $\mathbf2$ unequal to reflexivity.
This is an actual countercase to UIP for every type, not a semantic slogan.
\end{hbtheorem}

\begin{hbdefinition}{def:p07-hott-propositional-truncation}{The assumed propositional truncation interface}
For the following results, assume the source's additional type former
$\|A\|$, with $|a|:\|A\|$, specified paths between any two of its points,
and a recursor into a mere proposition $P$ from a map $A\to P$, computing
on $|a|$. Its existence is not derived from Id/J. Its higher-inductive
construction belongs to the separate HIT treatment.
The truncated disjunction is $\|P+Q\|$, and mere existential assertion is
$\|\Sigma_xP(x)\|$. An ordinary sum or dependent pair still retains data.
\end{hbdefinition}

\begin{hbtheorem}{thm:p07-hott-truncation-boundary}{proved=annex:p07-hott-truncation-boundary}{Unique choice and the obstruction to uniform witness extraction}
For a mere proposition $P$, $P\simeq\|P\|$. Consequently a family
of mere propositions $P(x)$ with $\Pi_x\|P(x)\|$ gives $\Pi_xP(x)$.
Under univalence and funext there is no uniform operation
$\Pi_{A:\mathcal U}(\|A\|\to A)$.
There is likewise no uniform double-negation elimination on all types,
nor an inhabitant of $\Pi_{A:\mathcal U}(A+(A\to\mathbf0))$.
This does not refute excluded middle restricted to mere propositions;
neither that axiom nor propositional resizing is assumed here.
\end{hbtheorem}

The untruncated dependent-pair distribution used in the proofs packages
already supplied witnesses; it is not the truncated choice axiom.
The latter cannot be inferred merely by moving a truncation through a
dependent product. The source distinguishes these tasks explicitly
\parencite[\S3.4--3.10]{HoTT2013PublicReleaseFoundations}.

\begin{hbtheorem}{thm:p07-hott-quasi-inverse-countercase}{proved=annex:p07-hott-quasi-inverse-countercase}{Why raw inverse data need not be a mere proposition}
Assuming funext,
$\mathsf{qinv}(\mathsf{id}_A)\simeq\Pi_{a:A}(a=a)$.
With univalence and propositional truncation, let
$X=\Sigma_{A:\mathcal U}\|\mathbf2=A\|$ in a larger universe.
Then $\mathsf{qinv}(\mathsf{id}_X)$ is not a mere proposition.
The proof constructs a nontrivial global loop family on this connected
component, not just one isolated nontrivial loop.
\end{hbtheorem}

\begin{proofboundary}{def:p07-hott-foundations-boundary}{Exact supplied scope and separate dependencies}
The constructions here give actual low-dimensional path operations,
transport, homotopy naturality, pair/function path comparisons,
equivalence-coherence/fiber bridges and the stated level closures.
They do not construct a full semantic infinity-groupoid model or all its
globular-operad coherences, prove univalence or funext from Id/J, or
establish canonicity for the added axioms. Higher-inductive constructions
continue in Chapter~\ref{sec:p07-higher-inductive-types}, and the distinct
computational cubical presentation is in Section~\ref{sec:p07-cchm2016}.
The higher homotopy-group, set-theoretic and categorical developments
require their own constructions. In particular, the natural-number code
family, algebraic-structure identity comparisons and remaining universal
properties are not consequences of the path laws alone.
\end{proofboundary}

\index{homotopy type theory}\index{identity type!path structure}
\index{transport!dependent}\index{equivalence!half adjoint}
\index{contractible type}\index{mere proposition}\index{homotopy n-type}
